% !PRINTED 590
\weberpage{590}
%<<CONT
the partial norm of $\mathfrak{p}$ is a power of $\mathfrak{P}$. We set
\weq{3}{\mathfrak{N}_R(\mathfrak{p}) = \mathfrak{P}^f,}
and call $f$ \gesperrt{the degree of $\mathfrak{p}$ with respect to the field $R$.}

The norm of $\mathfrak{P}$ in the field $R$ is likewise a power of $p$, which we shall denote by $P$; thus
\weq{4}{N_R(\mathfrak{P}) = P.}

For the total norm of $\mathfrak{P}$ in the field $\Omega$ we then obtain, by \S.\,159, (5) % [sic] source names $\mathfrak P$ here, while (5) and the norm identity require $\mathfrak p$
\weq{5}{N_\Omega(\mathfrak{p}) = P^f.}

If $\mathfrak{p}^g$ is the highest power of $\mathfrak{p}$ that divides $\mathfrak{P}$, then by 1. $\mathfrak{P}$ is also divisible by the $g$th power of all the prime factors conjugate to $\mathfrak{p}$, and by no higher power. Now if $\mathfrak{p}_1, \mathfrak{p}_2, \ldots, \mathfrak{p}_e$ are \gesperrt{the mutually distinct ones} among the prime ideals conjugate to $\mathfrak{p}$, then
\weq{6}{\mathfrak{P} = (\mathfrak{p}_1\mathfrak{p}_2\ldots\mathfrak{p}_e)^g,}
and if one takes the partial norm on both sides, and, as before, denotes by $n$ the degree of the field $\Omega$ with respect to $R$, so that the partial norm of $\mathfrak{P}$ becomes equal to $\mathfrak{P}^n$, then it follows from (3):
\weq{7}{n = efg.}
The natural numbers $e, f, g$ are therefore divisors of $n$.

If $\mathfrak{p}$ passes into $\mathfrak{p}_1$ under $\varphi_1$, that is, if
\[
\mathfrak{p}_1 = \mathfrak{p} \,|\, \varphi_1,
\]
then $\mathfrak{p} = \mathfrak{p}_1 \,|\, \varphi_1^{-1}$, and if $\Psi$ is the group associated with $\mathfrak{p}$, then for every substitution $\psi$ contained in $\Psi$ we also have
\[
\mathfrak{p}_1 = \mathfrak{p} \,|\, \psi \varphi_1, \quad \mathfrak{p}_1 = \mathfrak{p}_1 \,|\, \varphi_1^{-1} \psi \varphi_1.
\]

If, conversely, $\mathfrak{p}_1 \,|\, \varphi = \mathfrak{p}_1$, then it follows that $\mathfrak{p} \,|\, \varphi_1 \varphi \varphi_1^{-1} = \mathfrak{p}$, that is, $\varphi_1 \varphi \varphi_1^{-1} = \psi$ or $\varphi = \varphi_1^{-1} \psi \varphi_1$.

Hence $\mathfrak{p}_1$ is associated with the group $\varphi_1^{-1} \Psi \varphi_1$, and if
\[
\mathfrak{p}_1 = \mathfrak{p} \,|\, \varphi_1, \quad \mathfrak{p}_2 = \mathfrak{p} \,|\, \varphi_2, \ldots, \mathfrak{p}_e = \mathfrak{p} \,|\, \varphi_e
\]
then
\weq{8}{\Phi = \Psi\varphi_1 + \Psi\varphi_2 + \cdots + \Psi\varphi_e.}
$\Psi$ is thus a divisor of $\Phi$ of index $e$ and of degree $gf$.
